\documentclass[12pt,a4paper]{article}

\usepackage[T1]{fontenc}
\usepackage[utf8]{inputenc}
\usepackage[brazil]{babel}
\usepackage{amsmath,amssymb}
\usepackage{graphicx}
\usepackage{geometry}
\usepackage{booktabs}
\usepackage{longtable}
\usepackage{float}
\usepackage{hyperref}
\geometry{margin=2.5cm}
\hypersetup{hidelinks}

% Substitua os campos entre colchetes e insira as figuras indicadas.
% Compile com pdflatex relatorio.tex, preferencialmente duas vezes.
\newcommand{\pendente}[1]{\textbf{[PREENCHER: #1]}}

\begin{document}

\begin{titlepage}
\begin{center}
    \textbf{UNIVERSIDADE DE VILA VELHA}\\[0.3cm]
    \textbf{PROJETO E ANÁLISE DE ALGORITMOS}\\[4cm]
    {\Large\bfseries Explorando a intratabilidade e a força bruta no Problema do Caixeiro Viajante\par}
    \vspace{3cm}
    Bernardo Ziviani Gasparino\\
    Samuel Bruno da Silva\\
    Yuri Luiz Kruger Torquato\\[0.5cm]
    Professor: Wanderson Muniz de Santana\\[0.5cm]
    Turma: CC7N
    \vfill
    Vila Velha - ES\\
    2026
\end{center}
\end{titlepage}

\tableofcontents
\newpage

\section{Introdução}

O Problema do Caixeiro Viajante, conhecido pela sigla TSP, consiste em encontrar um circuito de custo mínimo que visite cada cidade uma única vez e retorne à cidade de origem. Embora sua definição seja simples, a quantidade de rotas candidatas cresce rapidamente com o número de cidades. Este trabalho investiga esse crescimento por meio de uma implementação de força bruta, relacionando as medições experimentais à análise apresentada no material da disciplina \cite{material} e aos objetivos da proposta de atividade \cite{proposta}.

O objetivo é implementar a busca exaustiva, registrar as quantidades de rotas e avaliações de arestas, medir o tempo para diferentes tamanhos de entrada e comparar os resultados com uma curva proporcional a $n!$.

\section{Fundamentação teórica}

\subsection{Modelagem e força bruta}

O conjunto de cidades é representado por $V=\{v_1,v_2,\ldots,v_n\}$, e o custo de deslocamento entre duas cidades é dado por $d(v_i,v_j)$. Na implementação, os custos são armazenados em uma matriz. Para cada rota candidata, o programa soma os custos das arestas percorridas e mantém a rota de menor custo encontrada.

Para evitar contar diferentes rotações do mesmo circuito, fixa-se $v_1$ como origem na execução principal. Restam, portanto, $n-1$ cidades a serem permutadas, resultando em
\begin{equation}
 R(n)=(n-1)!
\end{equation}
rotas. Cada rota possui $n$ arestas, incluindo o retorno à origem. O número total de avaliações dessas arestas é
\begin{equation}
 A(n)=n\cdot(n-1)!=n!.
\end{equation}
Assim, para a implementação sequencial com avaliação linear por rota, o tempo de execução é $\Theta(n!)$ \cite{material}. A igualdade acima representa o número de avaliações de arestas, não a quantidade total de instruções executadas pelo programa.

Como a matriz gerada neste projeto é simétrica, uma rota e sua inversa têm o mesmo custo. Mesmo assim, ambas são avaliadas separadamente, preservando a contagem $(n-1)!$ utilizada no material. Para $n=6$, por exemplo, são $5!=120$ rotas e $6\cdot120=720$ avaliações de arestas.

\subsection{Fixação da origem e redução do espaço de busca}

Sem fixar a origem, a enumeração de todas as permutações das $n$ cidades produziria $n!$ representações. Ao fixá-la, o total é reduzido para $(n-1)!$. Como $n!=n\cdot(n-1)!$, a redução é por um fator igual a $n$. As diferentes escolhas de origem correspondem a rotações de um mesmo circuito e não alteram seu custo ótimo.

O projeto também possui um modo de origem aleatória, que mantém a cidade sorteada fixa durante cada busca. Esse modo é complementar; para seguir literalmente a proposta, os resultados principais devem ser obtidos com origem fixa em $v_1$.

\section{Implementação}

\subsection{Algoritmo sequencial}

A implementação em Go gera uma matriz simétrica de custos positivos a partir de uma semente e enumera recursivamente as permutações das cidades restantes. Para cada permutação, constrói o circuito, avalia suas $n$ arestas e atualiza a melhor solução quando encontra um custo inferior. O programa registra a rota ótima, o custo, as contagens e o tempo de busca em arquivos CSV e de log com nomes únicos. O código-fonte está disponível no repositório do projeto \cite{repositorio}.

\section{Metodologia experimental}

\subsection{Ambiente e configuração}

Os experimentos devem identificar o ambiente e os parâmetros utilizados:
\begin{itemize}
    \item Processador e memória: \pendente{modelo e capacidade};
    \item Sistema operacional e versão do Go: \pendente{informações};
    \item Intervalo de cidades: \pendente{valores mínimo e máximo};
    \item Semente da matriz: \pendente{valor utilizado};
    \item Repetições por tamanho e tratamento das amostras: \pendente{informar procedimento}.
\end{itemize}

O tempo registrado compreende a enumeração, a montagem e a avaliação das rotas, incluindo a instrumentação. A geração da matriz e a escrita do CSV não fazem parte desse intervalo. O registro detalhado de todas as rotas não deve ser ativado nos testes de desempenho, pois a escrita no log pode afetar os tempos.

Cada execução do programa produz uma amostra por valor de $n$. Recomenda-se repetir as execuções e apresentar uma medida como a mediana, indicando o número de repetições. Os resultados e suas condições devem ser registrados sem substituir medições ausentes por estimativas.

\subsection{Comandos utilizados}

Para a execução principal com origem fixa, por exemplo:
\begin{verbatim}
make run ARGS="-min-n 3 -max-n 10"
\end{verbatim}
Para gerar o gráfico e a tabela a partir do CSV obtido:
\begin{verbatim}
make grafico RESULTADO="nome_do_arquivo.csv"
\end{verbatim}
O CSV é lido em \texttt{resultados/}, e o gráfico PNG e a tabela são gerados em \texttt{grafricos/}. Os nomes dos arquivos são únicos.

\section{Resultados e validação}

\subsection{Tabela de crescimento combinatório}

A Tabela~\ref{tab:resultados} deve ser preenchida com os dados efetivamente obtidos. A coluna $n!$ representa o número teórico de avaliações de arestas, enquanto as colunas de rotas e arestas medidas são extraídas do CSV. Acrescente ou remova linhas conforme o intervalo experimental.

\begin{longtable}{rrrrr}
\caption{Crescimento fatorial e resultados experimentais.}\label{tab:resultados}\\
\toprule
$n$ & $(n-1)!$ & $n!$ & Arestas medidas & Tempo (ms) \\
\midrule
\endfirsthead
\toprule
$n$ & $(n-1)!$ & $n!$ & Arestas medidas & Tempo (ms) \\
\midrule
\endhead
3 & 2 & 6 & \pendente{valor} & \pendente{valor} \\
4 & 6 & 24 & \pendente{valor} & \pendente{valor} \\
5 & 24 & 120 & \pendente{valor} & \pendente{valor} \\
6 & 120 & 720 & \pendente{valor} & \pendente{valor} \\
\bottomrule
\end{longtable}

\noindent\pendente{Adicionar as demais linhas até o maior valor de n testado e comparar também as rotas medidas com $(n-1)!$.}

\subsection{Curva teórica e dados empíricos}

O gráfico sobrepõe o tempo observado à função $f(n)=c\cdot n!$. A constante $c$ é ajustada pelo método dos mínimos quadrados aos tempos em milissegundos, considerando uma reta que passa pela origem. Para os pares $(x_i,y_i)=(n_i!,t_i)$, o valor utilizado é
\begin{equation}
 c=\frac{\sum_i x_i y_i}{\sum_i x_i^2}.
\end{equation}
O ajuste é realizado sobre os tempos em escala linear; consequentemente, as maiores medições têm maior influência. O gráfico utiliza escala logarítmica no eixo vertical apenas para facilitar a visualização de valores de ordens de grandeza diferentes. A proximidade entre as curvas ajuda na comparação, mas não constitui, isoladamente, uma prova da complexidade assintótica.

\begin{figure}[H]
    \centering
    % Substitua pelo caminho do gráfico PNG escolhido e remova o comentário:
    % \includegraphics[width=\linewidth]{grafricos/nome_do_grafico.png}
    \fbox{\parbox[c][5cm][c]{0.8\linewidth}{\centering Espaço para o gráfico de tempo medido e $c\cdot n!$}}
    \caption{Comparação entre o tempo medido e a curva fatorial calibrada.}
    \label{fig:grafico}
\end{figure}

\noindent\pendente{Informar o arquivo de dados, o valor de c, o intervalo testado e interpretar as diferenças observadas entre as curvas.}

\subsection{Evidências de execução}

\begin{figure}[H]
    \centering
    % Exemplo: \includegraphics[width=\linewidth]{caminho/print.png}
    \fbox{\parbox[c][4cm][c]{0.8\linewidth}{\centering Espaço para captura da execução ou dos resultados}}
    \caption{Evidência da execução do programa.}
\end{figure}

\noindent\pendente{Inserir capturas relevantes e descrever o que cada uma comprova.}

\section{Discussão das questões investigativas}

\subsection{Redução do espaço de busca}

A fixação de $v_1$ reduz a enumeração de $n!$ para $(n-1)!$ representações, um fator de redução igual a $n$. Essa economia deve ser verificada comparando a coluna de rotas avaliadas do CSV com $(n-1)!$ para cada entrada. \pendente{Apresentar um exemplo com os dados da execução.}

\subsection{Custo de avaliação por rota}

A avaliação percorre as $n$ arestas do circuito, inclusive a aresta de retorno. Dessa forma, o trabalho de avaliação de uma rota é $\Theta(n)$ e as avaliações de arestas nas $(n-1)!$ rotas somam $n!$. \pendente{Comparar a contagem registrada com a teoria para pelo menos um valor de n.}

\subsection{Limiar prático da intratabilidade}

O ponto em que a execução passa a levar minutos ou horas depende do hardware, da implementação e das condições de teste. \pendente{Informar o primeiro n observado com tempo em minutos ou horas, a duração, o arquivo de origem e o ambiente. Se o limiar não tiver sido atingido, declarar isso.}

O Merge Sort é um exemplo de algoritmo com complexidade $\Theta(n\log n)$, enquanto a busca exaustiva sequencial aqui estudada tem crescimento fatorial. A recorrência $T(n)=2T(n-1)+c$ discutida na seção 4.2 do material é um exemplo distinto de crescimento $\Theta(2^n)$; ela não é a recorrência usada para deduzir o custo deste TSP. A comparação evidencia que fatores polinomiais, exponenciais e fatoriais produzem comportamentos muito diferentes conforme a entrada aumenta.

\subsection{Classes P e NP e heurísticas}

A classe P reúne problemas de decisão solucionáveis em tempo polinomial; NP inclui problemas de decisão cujas soluções propostas podem ser verificadas em tempo polinomial. A versão de decisão do TSP é NP-completa, enquanto sua versão de otimização é NP-difícil. Não se conhece um algoritmo polinomial exato geral para o problema, mas isso não prova que tal algoritmo não exista: a questão $P=NP$ permanece em aberto.

Uma heurística gulosa pode escolher, em cada passo, uma cidade que pareça vantajosa localmente. Essa escolha não garante a melhor combinação global e, portanto, não assegura o custo mínimo para todas as instâncias. Heurísticas e metaheurísticas podem produzir soluções úteis em tempo menor, mas sua qualidade e suas garantias dependem do método empregado.

\section{Conclusão}

O estudo relacionou a enumeração de $(n-1)!$ rotas e a avaliação de $n$ arestas por rota à complexidade $\Theta(n!)$ da abordagem sequencial de força bruta. A fixação da origem elimina representações redundantes por rotação, mas não impede o crescimento fatorial.

\pendente{Sintetizar os resultados realmente obtidos, o maior n concluído, as limitações das medições e as conclusões extraídas do gráfico.}

\begin{thebibliography}{9}
\bibitem{material} UNIVERSIDADE DE VILA VELHA. \textit{Projeto e Análise de Algoritmos: Unidade 4}. Material da disciplina, seções 4.1 e 4.2.
\bibitem{proposta} \textit{Proposta de Atividade Investigativa: Explorando a Intratabilidade e a Força Bruta no Problema do Caixeiro Viajante (TSP)}. Documento fornecido para a atividade.
\bibitem{repositorio} CABRITTO CORPS. \textit{tsp-go: implementação do Problema do Caixeiro Viajante}. Repositório de código-fonte. Disponível em: \url{https://github.com/Cabritto-Corps/tsp-go}.
\end{thebibliography}

\end{document}
